% =====================================================================
\section{Arduino IDE и первая программа}
\label{les:ide}
% =====================================================================
\lessonintro
  {установить среду программирования, загрузить в плату первую программу, узнать правила записи программ и что такое переменные}
  {урок~\ref{les:board} (разъёмы и кнопки платы)}
  {плата присылает на компьютер приветствие и рассказывает о себе}
  {плата, USB-кабель \textbf{с передачей данных} (некоторые кабели умеют только заряжать!), компьютер}

\subsection{Как программа попадает в плату}

Микроконтроллер не понимает текст на человеческом языке. Мы пишем программу на языке C++, а
специальная программа --- \textbf{компилятор} --- переводит её в машинный код: длинную цепочку
нулей и единиц. Затем этот код по USB записывается во Flash-память платы. Это называется
\textbf{прошивкой}.

\begin{figure}[H]
\centering
\begin{tikzpicture}[node distance=6mm and 7mm]
  \node[blk,fill=blkmem] (src) {Программа\\ (скетч) \code{.ino}};
  \node[blk,fill=blkmem,right=of src] (cc) {Компилятор\\ \scriptsize переводит в\\ \scriptsize машинный код};
  \node[blk,fill=blkmem,right=of cc] (bin) {Файл\\ прошивки \code{.bin}};
  \node[blk,fill=blkper,right=of bin] (tool) {Загрузчик\\ esptool};
  \node[blk,fill=blkcpu,right=of tool] (chip) {Память\\ ESP32-S3};
  \draw[arr] (src) -- (cc);
  \draw[arr] (cc) -- (bin);
  \draw[arr] (bin) -- (tool);
  \draw[arr] (tool) -- node[above,font=\sffamily\scriptsize]{USB} (chip);
  \node[blk,fill=blkrf,below=8mm of cc] (lib) {Библиотеки\\ \scriptsize готовые функции};
  \draw[arr] (lib) -- (cc);
\end{tikzpicture}
\caption{Путь от текста программы до платы. Всё это Arduino IDE делает по одной кнопке}
\end{figure}

Программы для Arduino называют \textbf{скетчами} (англ. \emph{sketch} --- набросок).

\subsection{Установка Arduino IDE}

Мы будем писать программы в Arduino IDE 2.x --- бесплатной и простой среде.

\begin{enumerate}
  \item Скачай Arduino IDE 2 с сайта \url{https://www.arduino.cc/en/software} и установи.
  \item Открой \textsf{File → Preferences} (Файл → Настройки) и в поле
        \textsf{Additional boards manager URLs} вставь адрес:\\
    {\small\url{https://espressif.github.io/arduino-esp32/package_esp32_index.json}}
  \item Открой \textsf{Tools → Board → Boards Manager} (Инструменты → Плата → Менеджер плат),
        найди \textbf{esp32 by Espressif Systems} и установи версию 3.x. Скачивание займёт несколько минут.
  \item Подключи плату кабелем. В \textsf{Tools → Board} выбери \textbf{ESP32S3 Dev Module},
        а в \textsf{Tools → Port} --- появившийся порт (\code{COM3}, \code{/dev/ttyUSB0} или похожий).
  \item Если плата подключена через разъём \textbf{USB}, включи \textsf{Tools → USB CDC On Boot: Enabled}.
        Иначе сообщения от платы не дойдут до компьютера. Для разъёма UART эта настройка не нужна.
\end{enumerate}

\begin{tip}
Если порт не появился, попробуй другой кабель: очень часто виноват кабель «только для зарядки».
Для разъёма UART на Windows может понадобиться драйвер CP210x или CH343 --- его название обычно
написано на маленькой микросхеме рядом с разъёмом.
\end{tip}

\subsection{Из чего состоит скетч}

Любой скетч состоит из двух главных частей --- \textbf{функций}:
\begin{itemize}
  \item \code{setup()} (англ. «настройка») выполняется \emph{один раз} после включения. Здесь мы
        настраиваем выводы и готовимся к работе;
  \item \code{loop()} (англ. «петля») повторяется \emph{бесконечно}, пока есть питание. Здесь
        находится основная работа.
\end{itemize}

\begin{figure}[H]
\centering
\begin{tikzpicture}[node distance=6mm]
  \node[blk,fill=blkrtc] (rst) {Включение питания или кнопка RST};
  \node[blk,fill=blkmem,below=of rst] (boot) {Загрузчик запускает программу};
  \node[blk,fill=blkper,below=of boot] (setup) {\code{setup()} --- один раз};
  \node[blk,fill=blkper,below=of setup,minimum width=40mm] (loop) {\code{loop()} --- снова и снова};
  \draw[arr] (rst)--(boot); \draw[arr] (boot)--(setup); \draw[arr] (setup)--(loop);
  \draw[arr] (loop.east) -- ++(8mm,0) |- node[pos=0.25,right,font=\sffamily\scriptsize]{повторить} ($(loop.north east)+(0,1.5mm)$);
\end{tikzpicture}
\caption{Порядок работы программы}
\end{figure}

\subsection{Первая программа}

Первая программа ничего не подключает: она только сообщает о себе через USB. Набери её (лучше
именно набрать, а не скопировать --- так быстрее запоминается):

\codecap{Первый скетч: «Привет, ESP32-S3»}
\codeinput{l04_hello}

Что здесь написано:
\begin{itemize}
  \item \code{Serial.begin(115200)} --- включить связь с компьютером на скорости 115\,200 бит в секунду;
  \item \code{Serial.println("...")} --- отправить строку текста на компьютер;
  \item \code{Serial.printf(...)} --- отправить текст, подставив в него числа (вместо \code{\%lu});
  \item \code{delay(1000)} --- подождать 1000 миллисекунд, то есть одну секунду;
  \item \code{millis()} --- сколько миллисекунд прошло с момента включения.
\end{itemize}

Нажми кнопку \textsf{Upload} (стрелка вправо). Когда внизу появится \code{Done uploading}, открой
\textsf{Tools → Serial Monitor}, выбери скорость 115200 --- и увидишь сообщения от платы.
Как именно устроена эта связь, узнаем в уроке~\ref{les:serial}.

\subsubsection{Если прошивка не начинается}

Чтобы принять новую прошивку, чип должен войти в специальный режим. Обычно Arduino IDE делает это
сама, но если ты видишь ошибку \code{Failed to connect}, помоги ей кнопками:

\begin{center}
\begin{tikzpicture}[node distance=4mm]
  \node[blk,fill=blkper] (a) {Зажать\\ BOOT};
  \node[blk,fill=blkper,right=of a] (b) {Нажать и\\ отпустить RST};
  \node[blk,fill=blkper,right=of b] (c) {Отпустить\\ BOOT};
  \node[blk,fill=blkrf,right=of c] (d) {Нажать\\ «Upload»};
  \node[blk,fill=blkcpu,right=of d] (e) {После прошивки\\ нажать RST};
  \draw[arr] (a)--(b); \draw[arr] (b)--(c); \draw[arr] (c)--(d); \draw[arr] (d)--(e);
\end{tikzpicture}
\end{center}

\subsection{Учимся программировать: как записывается программа}

Программа похожа на рецепт: список команд, которые выполняются \emph{по порядку}, сверху вниз.
Компьютер выполняет их очень точно и ничего не додумывает. Если в рецепте написано «положить 100 кг
соли», компьютер положит 100 кг. Поэтому у языка C++ (читается «си-плюс-плюс»), на котором пишут
программы для Arduino, есть строгие правила записи. Посмотри на них на примере строчки из следующей
программы:

\begin{figure}[H]
\centering
\begin{tikzpicture}[tok/.style={font=\ttfamily\large,inner sep=1pt},ann/.style={font=\sffamily\scriptsize,align=center,text=black}]
  \node[tok,text=esp] (a) at (0,0) {int};
  \node[tok,text=accent,base right=0.5em of a] (b) {apples};
  \node[tok,base right=0.5em of b] (c) {=};
  \node[tok,text=okgreen!60!black,base right=0.5em of c] (d) {5};
  \node[tok,base right=0pt of d] (e) {\textbf{;}};
  \node[tok,text=gray,base right=1.2em of e] (f) {// яблок в корзине};
  \draw[esp,thick] (a.north) -- ++(0,0.35) node[ann,above] {тип:\\ что храним};
  \draw[accent,thick] (b.north) -- ++(0,1.1) node[ann,above] {имя\\ переменной};
  \draw[okgreen!60!black,thick] (d.north) -- ++(0,0.35) node[ann,above] {значение};
  \draw[black,thick] (e.south) -- ++(0,-0.35) node[ann,below] {точка с запятой ---\\ конец команды};
  \draw[gray,thick] (f.south) -- ++(0,-1.1) node[ann,below] {комментарий: для людей,\\ плата его не читает};
\end{tikzpicture}
\caption{Из чего состоит строчка программы}
\end{figure}

\begin{itemize}
  \item каждая команда заканчивается \textbf{точкой с запятой} \code{;};
  \item после \code{//} до конца строки идёт \textbf{комментарий} --- пояснение для людей;
  \item фигурные скобки \code{\{ \}} объединяют несколько команд в \textbf{блок}, как абзац в тексте.
        Всё, что между скобками после \code{void setup()}, --- это содержимое функции \code{setup()};
  \item большие и маленькие буквы различаются: \code{apples} и \code{Apples} --- разные имена;
  \item в дробных числах пишут \textbf{точку}, а не запятую: \code{23.5};
  \item текст пишут в двойных кавычках: \code{"Привет"}.
\end{itemize}

\begin{tip}
Если в программе ошибка, при компиляции внизу Arduino IDE появится красное сообщение с номером строки.
Чаще всего виноваты забытая \code{;}, незакрытая скобка или опечатка в имени. Ошибки --- это нормально:
их делают все программисты, даже самые опытные.
\end{tip}

\subsection{Учимся программировать: переменные}

Программе нужно что-то помнить: сколько раз нажали кнопку, какая яркость у светодиода, какая температура.
Для этого есть \textbf{переменные}. Переменная --- это коробка с названием, в которой хранится значение.
Значение можно положить, посмотреть и заменить другим. На коробке написано не только имя, но и
\textbf{тип} --- что в неё можно класть: целые числа, дробные, текст.

\begin{figure}[H]
\centering
\begin{tikzpicture}[
  box/.style={draw,thick,fill=#1,minimum width=24mm,minimum height=13mm,font=\ttfamily\Large},
  lbl/.style={font=\ttfamily\small,above=1mm},
  typ/.style={font=\sffamily\scriptsize,text=gray,below=1mm}]
  \node[box=blkmem] (a) at (0,0) {5};       \node[lbl] at (a.north) {apples};      \node[typ] at (a.south) {int --- целое};
  \node[box=blkper] (b) at (3.3,0) {23.5};  \node[lbl] at (b.north) {temperature}; \node[typ] at (b.south) {float --- дробное};
  \node[box=blkrf] (c) at (6.6,0) {true};   \node[lbl] at (c.north) {isOpen};      \node[typ] at (c.south) {bool --- да/нет};
  \node[box=blkrtc] (d) at (9.9,0) {"Маша"};\node[lbl] at (d.north) {name};        \node[typ] at (d.south) {String --- текст};
\end{tikzpicture}
\caption{Переменные --- подписанные коробки. Тип говорит, что в коробку можно положить}
\end{figure}

\begin{table}[H]
\centering\small
\begin{tabularx}{\textwidth}{@{}l X l@{}}
\toprule
Тип & Что хранит & Пример\\
\midrule
\code{int} & целое число от $-2\,147\,483\,648$ до $2\,147\,483\,647$ (32 бита) & \code{int score = 10;}\\
\code{float} & дробное число (около 7 точных цифр) & \code{float t = 36.6;}\\
\code{bool} & логическое значение: \code{true} (истина) или \code{false} (ложь) & \code{bool on = false;}\\
\code{String} & строку текста & \code{String s = "Привет";}\\
\code{uint8\_t}, \code{byte} & целое от 0 до 255 (8 бит, как в уроке~\ref{les:signals}) & \code{byte r = 200;}\\
\code{uint32\_t} & целое от 0 до 4\,294\,967\,295 --- для времени \code{millis()} & \code{uint32\_t t = millis();}\\
\bottomrule
\end{tabularx}
\caption{Типы переменных, которые встретятся в курсе}
\end{table}

Знак \code{=} в программировании означает не «равно», а «\textbf{положить в коробку}». Запись
\code{apples = apples + 3} читается так: «возьми, что лежит в \code{apples}, прибавь 3 и положи результат
обратно в \code{apples}». В математике это бессмыслица, а для программы --- обычное дело.

\begin{figure}[H]
\centering
\begin{tikzpicture}[box/.style={draw,thick,fill=blkmem,minimum width=16mm,minimum height=11mm,font=\ttfamily\Large}]
  \node[box] (b1) at (0,0) {5};
  \node[font=\ttfamily\small,above=1mm of b1] {apples};
  \node[font=\sffamily\small,below=1mm of b1] {было};
  \node[draw,thick,rounded corners,fill=blkper,font=\ttfamily\small,align=center] (op) at (4,0) {5 + 3 = 8};
  \node[box,fill=okgreen!20] (b2) at (8,0) {8};
  \node[font=\ttfamily\small,above=1mm of b2] {apples};
  \node[font=\sffamily\small,below=1mm of b2] {стало};
  \draw[arr] (b1) -- node[above,font=\sffamily\scriptsize]{достать} (op);
  \draw[arr] (op) -- node[above,font=\sffamily\scriptsize]{положить} (b2);
  \node[font=\ttfamily\large] at (4,1.4) {apples = apples + 3;};
\end{tikzpicture}
\caption{Присваивание: сначала считается правая часть, потом результат кладётся в переменную}
\end{figure}

\begin{table}[H]
\centering\small
\begin{tabular}{@{}lll@{}}
\toprule
Запись & Что делает & Пример ($x = 7$)\\
\midrule
\code{+ - * /} & сложение, вычитание, умножение, деление & \code{x * 2} → 14\\
\code{\%} & остаток от деления & \code{x \% 2} → 1\\
\code{x++}, \code{x--} & увеличить или уменьшить на 1 & \code{x} станет 8 или 6\\
\code{x += 5} & то же, что \code{x = x + 5} & \code{x} станет 12\\
\bottomrule
\end{tabular}
\caption{Арифметические операции}
\end{table}

\codecap{Переменные и операции}
\codeinput{l04_variables}

\begin{caution}
Если оба числа целые, деление тоже будет целым: \code{7 / 2} даёт \textbf{3}, а не 3,5 --- дробная часть
просто отбрасывается. Чтобы получить 3,5, хотя бы одно число должно быть дробным: \code{7.0 / 2}.
\end{caution}

Переменная, объявленная \emph{вне} функций (в начале программы), --- \textbf{глобальная}: её видят все части
программы. Переменная, объявленная \emph{внутри} фигурных скобок, --- \textbf{локальная}: она существует только
внутри этого блока. Как записка на доске в классе (видят все) и записка в твоём кармане (видишь только ты).

Теперь понятно, что было в первой программе: \code{setup()} и \code{loop()} --- это \textbf{функции}, наборы
команд под одним именем; \code{Serial.println()} и \code{delay()} --- готовые функции, которые мы
\emph{вызываем}, передавая им в скобках данные; \code{millis()} возвращает число --- его можно положить в переменную.

\begin{summary}
\begin{itemize}[leftmargin=*]
  \item Программа (скетч) пишется в Arduino IDE, переводится компилятором в машинный код и загружается по USB.
  \item \code{setup()} выполняется один раз, \code{loop()} --- бесконечно.
  \item \code{Serial} позволяет плате отправлять сообщения на компьютер.
  \item Если прошивка не идёт --- кнопки BOOT и RST.
  \item Команды заканчиваются \code{;}, блоки --- в \code{\{ \}}, комментарии --- после \code{//}.
  \item Переменная --- коробка с именем и типом; \code{=} кладёт в неё значение; целые числа делятся нацело.
\end{itemize}
\end{summary}

\begin{tasks}
\begin{enumerate}
  \item Добавь в программу вывод своего имени и возраста.
  \item Сделай так, чтобы плата печатала время работы не в миллисекундах, а в секундах.
  \item Узнай, сколько памяти (Flash и PSRAM) у твоей платы, и сравни с надписью на модуле.
  \item Заведи переменные \code{a = 17} и \code{b = 5} и напечатай их сумму, разность, произведение, частное и остаток.
        Почему частное получилось таким?
\end{enumerate}
\end{tasks}
